\subsection{Agent Governance Landscape}\label{subsec:agent-governance}

LLM agent proliferation has created an urgent need for governance beyond traditional ontology engineering. ADL Lite fills a gap that trust, safety, and provenance layers overlook: a lightweight, event-sourced, tamper-evident registry of agent capabilities and their lifecycle states.

\textbf{AgentHub}~\cite{agenthub2025} proposes a registry layer for agent sharing with discovery, workflow integration, and lifecycle transparency, but it remains a research agenda without an implemented lifecycle state machine or deterministic derivation semantics. ADL Lite extends beyond AgentHub's agenda by providing an implemented registry with formal lifecycle semantics (register $ \to $ validate $ \to $ deprecate), cryptographic hash chaining, and deterministic status derivation.\par
\textbf{KYA}~\cite{kya2025} verifies agent identity via HMAC chains but does not audit capability lifecycles. \textbf{Zhou}~\cite{zhou2026governance} proposes cryptographic binding and reproducibility verification for agent tool use, addressing tool-call provenance at the invocation level; ADL Lite audits the capability lifecycle that such tool use presupposes. \textbf{AgentSafe}~\cite{agentsafe2025} offers comprehensive architecture governance but requires heavyweight infrastructure and lacks Markdown-native authoring. \textbf{From Agent Traces to Trust}~\cite{agenttraces2025} concludes that ``a unified trace schema is still needed''---ADL Lite provides a unified capability-lifecycle trace schema. \textbf{Talukdar et al.}~\cite{talukdar2025} demonstrate multi-agent ontology engineering but produce static artifacts without lifecycle audit. \textbf{SafeAgent}~\cite{safeagent2025} operates at the interaction layer rather than the artifact layer. Runtime safety systems (SafeHarness~\cite{safeharness2026}, OpenClaw PRISM~\cite{openclaw_prism2026}) and runtime infrastructure (Institutional AI~\cite{institutional_ai2026}, GaaS~\cite{gaas2025}, COMPASS~\cite{compass2026}, UALM~\cite{ualm2026}) consume capability provenance for policy decisions; ADL Lite occupies the \emph{registry layer} between permissioning and runtime enforcement.

Table~\ref{tab:agent-governance} compares these systems.

\begin{table}[htbp]
\centering
\footnotesize
\caption{Comparison of agent governance systems.}
\label{tab:agent-governance}
\begin{tabularx}{\textwidth}{@{}p{2.2cm}XXXXXX@{}}
\toprule
\textbf{Dimension} & \textbf{KYA} & \textbf{AgentSafe} & \textbf{Agent Traces} & \textbf{Talukdar et al.} & \textbf{SafeAgent} & \textbf{ADL Lite} \\
\midrule
Governance Scope & Trust / permissions & Architecture safety & Provenance / trace & Ontology production & Protocol safety & Capability lifecycle \\
Lifecycle Awareness & None & Design-time only & None & None & None & Full (register $\to$\allowbreak validate $\to$\allowbreak deprecate) \\
Deployment Weight & Lightweight & Heavyweight & Survey / framework & Research pipeline & Lightweight & Very lightweight (pip + Git) \\
Tamper-evidence & HMAC chain & Audit logs & Trace schema (proposed) & Manual QA review & Sandbox checks & SHA-256 hash chain + preconditions \\
Agent-Native & Yes (15+ adapters) & Yes (design-time) & Yes (provenance) & Yes (multi-agent) & Yes (runtime) & Yes (Markdown-native, multi-agent) \\
\bottomrule
\end{tabularx}
\end{table}

Runtime safety systems (SafeHarness~\cite{safeharness2026}, OpenClaw PRISM~\cite{openclaw_prism2026}) govern execution-time safety rather than capability registration; they complement ADL Lite by leveraging the capability metadata that ADL Lite produces.

\textbf{Agent capability standards.} Emerging agent standards---the Model Context Protocol (MCP)~\cite{mcp2024anthropic}, Google's Agent2Agent (A2A)~\cite{a2a2025google}, and the Linux Foundation's AGNTCY Open Agent Schema Framework~\cite{agntcy2025linux}---provide agent capability registration and discovery, but they do not provide lifecycle audit semantics: none offers an append-only, tamper-evident chain of capability states, LUB-derived status, or precondition enforcement over lifecycle transitions. General-purpose orchestration frameworks (LangChain~\cite{langchain}, LlamaIndex~\cite{llamaindex}, the OpenAI Agents SDK~\cite{openai_agents}) supply tool-calling and workflow primitives but likewise treat capability metadata as ephemeral runtime state rather than an auditable lifecycle. This is precisely the layer that ADL Lite adds on top of such standards and frameworks (Section~\ref{subsec:complementary-systems}).

\subsection{Foundational Ontologies and Event-First Semantics}\label{subsec:foundational-ontologies}

ADL Lite's \texttt{Event} corresponds to BFO's \emph{process}~\cite{bfo}, DOLCE's \emph{perdurant}~\cite{dolce}, and the event/action distinction in UFO~\cite{ufo}. ADL Lite adopts \emph{pragmatic pluralism}: BFO as conceptual guide, DOLCE as design principle, and UFO as modelling framework. A detailed cross-mapping appears in Section~\ref{sec:ontological-analysis} (Table~\ref{tab:upper-ontology-mapping}).\par
We next review knowledge graph construction and ontology engineering systems that provide the broader tooling context for ADL Lite's document-native approach.

\subsection{Knowledge Graph Construction and Ontology Engineering}\label{subsec:kg-ontology-engineering}

The Semantic Web vision~\cite{bernerslee2001semanticweb} was operationalised through OWL~2~\cite{owl2}; toolchains enable ontology authoring at scale, but reasoners impose computational costs that render them ill-suited to rapidly evolving knowledge bases~\cite{glimm2014hermit,sirin2007pellet,hogan2021knowledge}. The OBO Foundry provides centralised versioning~\cite{obo_foundry}; lightweight alternatives (Schema.org~\cite{guha2016schemaorg}, JSON-LD~\cite{json_ld}, SHACL~\cite{w3c_shacl}, ShEx~\cite{shex}) and Markdown tools~\cite{foam2020,dendron2020,obsidian2024} lower the barrier to entry but either remain bound to RDF or lack machine-checkable coordination. Industry-scale knowledge graphs~\cite{noy2019industrykg} and collaborative knowledge bases such as Wikidata~\cite{wikidata_vrandecic2014} show the curation challenge at scale. A systematic literature review of LLM-based ontology engineering~\cite{garijo2025llmoe} catalogues this landscape; LLM4ACOE~\cite{llm4acoe}, Sim-HCOME~\cite{sim_hcome}, and Ontology-Constrained Neural Reasoning~\cite{tuan2026ontology} generate OWL automatically but without lifecycle audit. ADL Lite differs from these systems in three respects: (i)~\emph{authoring paradigm:} ADL Lite uses Markdown-native documents rather than OWL/RDF output; (ii)~\emph{audit scope:} ADL Lite provides built-in lifecycle state machines (register $ \to $ validate $ \to $ deprecate) rather than static ontology production; (iii)~\emph{deployment weight:} ADL Lite is pip-installable and operates on Git repositories, whereas LLM4ACOE and Sim-HCOME require multi-agent orchestration pipelines. These are complementary trade-offs, not competitive disadvantages: LLM4ACOE/Sim-HCOME excel at automated ontology generation with ReAct reasoning traces, while ADL Lite excels at lightweight, auditable management of existing capabilities. Table~\ref{tab:llm-oe-comparison} presents the dimension-level comparison.

\begin{table}[htbp]
\centering
\footnotesize
\caption{Dimension-level comparison: ADL Lite vs. proximate systems.}
\label{tab:llm-oe-comparison}
\begin{tabularx}{\textwidth}{@{}p{3.2cm}XXX@{}}
\toprule
\textbf{Dimension} & \textbf{LLM4ACOE / Sim-HCOME} & \textbf{Blocklace (2024)} & \textbf{ADL Lite} \\
\midrule
Primary output & OWL~2~DL ontology & Generic BFT-CRDT DAG & Markdown-native capability registry \\
LLM agent role & Multi-agent ReAct ontology generation & None (transport layer) & Actor in lifecycle events \\
Lifecycle audit & None (static ontology) & None (application-agnostic) & Register $\to$\allowbreak validate $\to$\allowbreak deprecate \\
Cryptographic integrity & None (manual QA review) & SHA-256 + BFT consensus & SHA-256 hash chain \\
Deterministic state & External reasoner required & None (raw data structure) & Built-in $\delta(C)$, $\gamma(C)$ \\
Deployment & Research pipeline (multi-agent) & Library (distributed systems) & \texttt{pip install} + Git \\
\bottomrule
\end{tabularx}
\end{table}

\textbf{Comparison with FAOS.} The Foundation AgenticOS (FAOS) platform \cite{tuan2026ontology} proposes a three-layer enterprise ontology (Role, Domain, Interaction) for grounding LLM agents in organisational contexts. Whereas FAOS's layers govern \emph{agent roles} and \emph{domain interactions} within an enterprise, ADL Lite's L1--L4 layers audit \emph{capability lifecycle events} across decentralized multi-agent ecosystems; FAOS does not provide cryptographic hash chaining, deterministic status derivation, or append-only event sourcing. The two are complementary: ADL Lite could consume FAOS-generated ontologies and audit their lifecycle through EventChains.

Blocklace~\cite{blocklace2024} is a general-purpose BFT-CRDT with a cryptographic hash DAG, providing Byzantine fault tolerance and equivocation detection at the transport layer. ADL Lite is complementary: Blocklace provides the \emph{transport} (distributed consensus over a DAG of events), while ADL Lite provides the \emph{application semantics} (deterministic status derivation $\delta(C)$, confidence aggregation $\gamma(C)$, and declarative preconditions); the proposed Phase~3 architecture uses Blocklace's DAG beneath ADL Lite's lifecycle semantics. SEO~2025~\cite{boldachev2025seo} uses causal DAGs with reactive guards for distributed multi-agent event contribution; ADL Lite's linear hash chain is less expressive but offers $O(1)$ per-event append latency and deterministic $O(n)$ verification---a deliberate expressivity--simplicity trade-off.

\subsection{Provenance, Trust, and Verifiable Publishing}\label{subsec:provenance-trust}

W3C PROV-O~\cite{w3c_prov_o} is the W3C standard provenance interchange model, surveyed in detail by Moreau and Groth~\cite{moreau2013prov}; the earlier named-graph account~\cite{carroll2005namedgraphs} already identified graph provenance as a basis for trust on the Web. Linked Data Proofs~\cite{linked_data_proofs} and RDF-star~\cite{rdf_star} provide statement-level annotation but lack both cryptographic integrity and lifecycle audit. Secure-provenance research~\cite{braun2010secure,hasan2009preventing} formalises adversarial tamper-evidence for provenance chains but targets low-level data lifecycle management rather than capability lifecycle audit. Decentralised identifiers and verifiable credentials (DID/VC)~\cite{mazzocca2024didvc} provide strong actor authentication but weak content binding: a DID identifies \emph{who} made a claim, not \emph{what} the claim is about. ADL Lite complements DID-based identity by binding the genesis hash to the specific event (Section~\ref{subsec:identity-resilience}). Git-native systems (RO-Crate~\cite{ro_crate}, DataLad~\cite{datalad}, TerminusDB~\cite{terminusdb}) share immutable history but lack document-native authoring. Full mappings to standards and loss analysis are in Appendix~F.

\subsection{Event-Centric Ontologies}\label{subsec:event-centric}

CIDOC CRM~\cite{cidoc_crm}, SEM/LODE~\cite{vanhage2011sem,shaw2011lode}, and RDF Stream Processing~\cite{w3c_rsp,c_sparql,lephuoc2011cqels} provide event vocabularies but lack both constraint mechanisms and retrospective verification. Boldachev's SEO~\cite{boldachev2025seo}, Al-Fedaghi's occurrence-only paradigm~\cite{alfedaghi2024occurrence}, and Guizzardi's UFO-B~\cite{guizzardi2024ufob} offer alternative event models. CRDTs~\cite{crdts,automerge,yjs,crdt_rdf_2023,nieto2024crdt} and Blocklace~\cite{blocklace2024} reconcile concurrent updates; ADL Lite complements these with append-only event chains and deterministic lifecycle audit. Real-time collaborative ontology editing using distributed version control (OntoEditor, which uses Operational Transformation rather than CRDTs) provides a related but distinct approach to multi-author ontology development~\cite{hemid2024ontoeditor}. ADL Lite's EventChain extends the event-sourced aggregate root pattern~\cite{event_sourcing,event_sourcing_fowler} with ontological analysis, SHA-256 chaining, and Markdown-native authoring, while its separation of append-only event streams from derived query state follows the command-query responsibility segregation (CQRS) pattern~\cite{cqrs}.

\subsection{Transparency Logs and Software Supply Chain}\label{subsec:transparency-logs}

Certificate Transparency~(CT)~\cite{rfc6962,crosby_wallach}, Sigstore~Rekor~\cite{sigstore2022}, in-toto~\cite{in_toto}, and SLSA~\cite{slsa} provide tamper-evident registries for software artifacts, certificates, and supply-chain provenance. These systems use Merkle-tree or hash-DAG structures to enable $O(\log n)$ inclusion proofs and $O(1)$ consistency proofs, making them highly scalable for globally distributed verification with millions of entries. ADL Lite's linear SHA-256 chain, by contrast, requires $O(n)$ sequential verification for full integrity checking---a deliberate design trade-off: linear chains are simpler to implement in Markdown-native environments, produce human-readable audit trails, and require no external gossiping or witness infrastructure. Merkle trees excel at proving that an entry exists in a large append-only log; ADL Lite excels at auditing the lifecycle of a single capability through deterministic status derivation and confidence aggregation. The two are complementary: future Phase~3 integration (FW4, FW9) may adopt Merkle-tree batching for cross-repository verification while retaining per-chain lifecycle semantics. Blockchain-based provenance~\cite{akbarfam2024blockchain} and the records-management lifecycle view of distributed ledgers~\cite{barbereau2026lifecycle} reinforce this design: Barbereau et al.\ reconceptualize blockchain systems as record-management systems with a seven-stage lifecycle (ISO 15489-1:2016), echoing the register $\to$ validate $\to$ deprecate lifecycle that ADL Lite makes native and deterministic.\par
\textbf{Quantitative comparison with nanopublications and Git-native provenance.} Supplementary Table S19 compares ADL Lite with two related lightweight approaches on operational metrics relevant to capability-lifecycle audit.

Nanopublications with trusty URIs~\cite{nanopub2023} provide $O(1)$ hash verification for individual assertions but lack lifecycle semantics: each nanopublication is a standalone, immutable statement with no formal status derivation or precondition-governed transitions. Recent NanoPub-Jelly streaming formats (2024) support chained nanopublications with temporal ordering, though without the lifecycle audit semantics (register $ \to $ validate $ \to $ deprecate) and deterministic status derivation ($\delta$, $\gamma$) that ADL Lite provides. Git-native systems (RO-Crate, DataLad) provide commit-level provenance but treat the entire commit as a unit of change, making it impossible to trace the lifecycle of an individual capability within a repository containing many concepts. ADL Lite trades $O(1)$ verification for $O(n)$ chain scanning in exchange for per-event lifecycle audit: each event (REGISTER, VALIDATE, DEPRECATE, FORK) is individually typed, preconditioned, and auditable. The migration path to nanopublications is bidirectional: ADL Lite events export to RDF-star quoted triples (Appendix~F), and nanopublication assertions can be imported as RELATE events with SHA-256 content-addressed references. The migration path to Git-native provenance is direct: ADL Lite EventChains are stored as Markdown files in standard Git repositories, with commit signatures providing repository-level integrity and event hashes providing concept-level integrity.

\subsection{Deep Comparison with Nanopublications}\label{subsec:nanopub-comparison}

Nanopublications with trusty URIs~\cite{nanopub2023} are the closest Semantic Web analogue to ADL Lite's per-event integrity and provenance tracking; Supplementary Table S21 summarises the dimension-level comparison.

\textbf{Immutability granularity: per-capability vs. per-statement.} Nanopublications enforce immutability at the \emph{statement} level: each is a single signed RDF assertion with a trusty URI containing a SHA-256 hash; retraction is a new nanopublication linked via \texttt{prov:wasDerivedFrom}. ADL Lite enforces immutability at the \emph{capability} level: an EventChain is an append-only sequence of typed events that collectively govern a single concept's lifecycle, and the chain's state is derived from the entire history---appropriate for capability governance, where the \emph{process} of validation, deprecation, and forking matters as much as any individual assertion.

\textbf{Key structure: trusty URI vs. genesis hash.} A nanopublication trusty URI encodes the publisher's RSA key, creation timestamp, and the content hash; the URI is immutable, so renaming requires publishing a new nanopublication. In ADL Lite, the \texttt{concept\_id} is a human-readable alias while the genesis hash (SHA-256 of the first event) is the operational identifier---a content-derived property, not part of the URI---allowing concepts to be renamed (via a fork) while preserving the cryptographic identity of the original chain.

\textbf{Lifecycle semantics: absent vs. native.} Nanopublications have no built-in lifecycle semantics: an assertion's status is a consumer inference, not encoded in the format. ADL Lite natively encodes lifecycle events (\texttt{REGISTER}, \texttt{VALIDATE}, \texttt{DEPRECATE}, \texttt{FORK}, \texttt{ARCHIVE}) as first-class chain entities with deterministic status derivation $\delta(C)$ and confidence aggregation $\gamma(C)$; the status of a concept is a mathematically defined function of the chain's history, not a consumer inference.

\textbf{Worked example: capability deprecation.} Deprecating a validated capability in a nanopublication pipeline requires publishing a new ``deprecated'' nanopublication, linking it via \texttt{prov:wasDerivedFrom}, and relying on consumers to query both and infer supersession---there is no mechanism to enforce that the link is created or interpreted correctly. In ADL Lite, deprecation is a single \texttt{DEPRECATE} event appended to the existing chain: typed, preconditioned (requires current status \texttt{validated}), hash-linked, and immediately making $\delta(C) = \texttt{deprecated}$ for all consumers, with no ambiguity and a complete audit trail.

\textbf{Export/import mapping.} The migration path is bidirectional. ADL Lite exports to nanopublication-like structures via RDF-star---each event becomes a quoted triple with its event hash as provenance, e.g.:
\begin{lstlisting}[language=Turtle]
<< :gradient-explosion :isomorphic-to :vanishing-gradient >>
    prov:wasGeneratedBy :evt-gradient-explosion-validate-001 ;
    adl:eventHash "sha256:a3f2..." ;
    adl:confidence 0.91 .
\end{lstlisting}
Conversely, a nanopublication assertion can be imported as an ADL Lite \texttt{RELATE} event with the trusty URI hash stored as a content-addressed reference, enabling ADL Lite to act as a lifecycle audit layer \emph{on top of} existing nanopublication repositories.

\subsection{Ontology Registry Comparison}\label{subsec:registry-comparison}

ADL Lite occupies a distinct position relative to mature ontology and semantic-artifact registries that provide curation, versioning, and governance for ontologies and vocabularies (Supplementary Table S22).

BioPortal and OntoPortal \cite{bioportal,ontoportal} are the de facto infrastructure for ontology hosting in the biomedical domain. They offer versioning via release tags, collaborative submission, and Web-based curation workflows, but they do not record the \emph{process} of validation: a submitter uploads an OWL file, a curator reviews it, and the ontology is published---the curation \emph{decision} is not itself an event in the ontology's lifecycle. FAIRsharing \cite{fairsharing} curates standards and databases but focuses on metadata discovery rather than lifecycle audit. ADL Lite complements these registries by providing (i) event-sourced lifecycle audit, in which each validation, deprecation, and fork is a hash-linked event; (ii) multi-agent native authoring with deterministic derivation semantics ($\delta$, $\gamma$); and (iii) content-addressability alongside provenance-based identity, enabling automated duplicate detection. The migration path is bidirectional: ADL Lite concepts export to OWL 2 DL (Appendix~\ref{app:owl-dl}), and existing BioPortal ontologies can be imported as initial \textsc{Register} events with their URI preserved as a \texttt{RELATE} target.

\textbf{Release-based versus event-sourced versioning.} Ontology versioning on the Semantic Web has long been recognised as a hard problem~\cite{klein_fensel_2001}; BioPortal and OntoPortal use \emph{release-based} versioning: an ontology is uploaded as a monolithic file, and each release is a snapshot with a version tag (e.g., ``v2024-01''). This model captures the \emph{state} of the ontology at discrete points but does not capture the \emph{process} by which that state was reached: who proposed a change, who validated it, and what evidence was provided. ADL Lite's \emph{event-sourced} model treats every lifecycle action (registration, validation, deprecation, fork) as an atomic event, enabling fine-grained provenance and deterministic state reconstruction from the event history. Release-based versioning is appropriate for mature, slowly evolving ontologies; event-sourced versioning is necessary for rapidly evolving, multi-agent capability ecosystems where decisions must be auditable and reversible.

\subsection{Positioning}\label{subsec:related-positioning}

Palantir Foundry unifies ``semantic elements'' (nouns) with ``kinetic elements'' (verbs)~\cite{palantir_fde}, but requires enterprise-scale deployment and proprietary tooling, and lacks LLM-native authorship. ADL Lite offers pip-installable deployment, Markdown-native authoring, and multi-agent deterministic derivation. Full comparison tables are in Appendix~F.

Within the open-source, pip-installable, Markdown-native tool space, to the best of our knowledge, ADL Lite uniquely combines four properties: (a)~document-native authoring in plain Markdown, (b)~built-in tamper detection through SHA-256 hash chaining, (c)~lifecycle state machines with declarative preconditions, and (d)~pip-installable deployment requiring no enterprise infrastructure. This enables multiple LLM agents to propose, challenge, and refine concepts within a shared document-native ontology, with the EventChain providing the audit trail and deterministic derivation mechanism that renders decentralized authorship accountable.
